\section{Introduction}
\label{sec:intro}

Enterprises record regulations, failure analyses and operational lessons in
unstructured documents. Answering a question such as \emph{why did edge
delamination appear in a photovoltaic inspection?} requires finding the report
that records the decision, reading its evidence and citing it accurately.
Retrieval-augmented generation (RAG) addresses document discovery by embedding
queries and passages in a shared vector space \cite{rag-lewis,dpr}; agentic
systems extend this process by letting a language model call retrieval tools
\cite{mcp,singh-agentic-rag}. Yet semantic similarity measures topical
relatedness, not whether a passage contains evidence for the requested claim
\cite{siren,gao-ragsurvey}.

The relevant capabilities are usually treated separately. Retrieval asks
whether a relevant passage enters a ranked list. Knowledge management governs
how documents are organized and maintained. Document AI recovers text, tables
and formulas from complex layouts \cite{mineru}. In deployed systems, these
operations form one dependency chain: parsing errors affect indexing, placement
limits what retrieval can see, and ranking alone cannot establish whether a
candidate supports an answer. No single stage owns the property the user
actually needs---that the sentence being cited is supported by the document
being cited---and so no stage is accountable for it.

This paper presents QDCVR (Query-Driven Content-Verified Retrieval), a
self-hosted knowledge-base management system that makes content the governing
object across acquisition, organization and question answering, and assigns
that accountability to three enforceable contracts:

\begin{enumerate}
\item \emph{Content governs placement.} The parsed body, not the filename or
the source directory, decides which library holds a document, what its
description and tags say, and therefore what a later search is able to see.
\item \emph{Content governs admission.} A retrieved candidate becomes
admissible evidence only after the server fetches its text and a content gate
scores it; rank position confers no authority of its own.
\item \emph{Absence is reported.} When no fetched segment meets the evidence
criterion, the system returns an explicit not-found report naming the scope it
searched, instead of a plausible unsupported answer drawn from parametric
knowledge.
\end{enumerate}

Enforcement is server-side --- quality checks, lifecycle audits, lane-specific
tool policies and a fail-closed evidence gate --- and each contract is measured
on a deployed collection rather than asserted from the implementation. The
deployed platform integrates document parsing, content-routed organization,
lifecycle management, operational experiences, multi-lane retrieval and a Model
Context Protocol (MCP) server with 95 registered tools \cite{mcp}, runs
locally, and exposes the same retrieval functions to several agent harnesses;
untested functions are identified rather than inferred from their
implementation.

\noindent\textbf{Contributions.}
\begin{enumerate}
\item A content-governance architecture whose three contracts are enforced by
mechanism rather than convention: quality-gated acquisition, content-derived
placement, three-source consistency checks and evidence-based answer admission
(Sections~\ref{sec:system}--\ref{sec:implementation}).
\item A verified-retrieval design that separates candidate discovery from
evidence admission: three budgeted lanes over one server-side content gate,
with an explicit abstention contract. The evaluation attributes outcomes to the
verdict's role on identical candidate pools---the same verdicts work as a
filter and fail as a re-ranker (Section~\ref{sec:lanes} and
Section~\ref{sec:e4}).
\item A deployment-based evaluation spanning three corpora that reports
negative results alongside successful checks: papers that never became
retrievable and an inventory that cannot attribute them, a gate with no
measured filtering power on the industrial corpus, and ablations that return
null (Section~\ref{sec:eval}).
\end{enumerate}

A CIKM~2026 demonstration introduced the organization and retrieval layers
\cite{qdcvr-demo}. This article adds the acquisition pipeline, lifecycle
contracts, operational-knowledge store and a broader deployment evaluation.

\subsection{Research questions and reading guide}
\label{sec:rq}

Five questions structure the design-to-evidence sequence: how content
commitments can be enforced across retrieval strategies and agent interfaces
(RQ1); how parsing gates, content routing and lifecycle checks behave on
deployed collections (RQ2); how the retrieval lanes balance document recall,
evidence verification and operating cost (RQ3); under what conditions content
verification helps, fails or requires calibration (RQ4); and where abstention,
latency, cost and reliability constraints arise (RQ5).

Section~\ref{sec:related} positions the work; Sections~\ref{sec:system}--
\ref{sec:organization} describe the architecture and management functions;
Section~\ref{sec:lanes} specifies retrieval and verification;
Section~\ref{sec:implementation} covers deployment and enforcement; and
Section~\ref{sec:eval} evaluates the system, followed by implications,
limitations and conclusions in
Sections~\ref{sec:discussion}--\ref{sec:conclusion}.
